Pendant un orage, Léa voit un éclair, puis entend le tonnerre \qty{\v{t}}{s} plus tard. La vitesse du son dans l'air est de \qty{\v{v}}{m/s} et celle de la lumière est de \qty{3.00e8}{m/s}.

\begin{enumerate}[label=\alph*)]
    \item À quelle distance de Léa l'éclair s'est-il produit ? (On peut considérer que la lumière arrive instantanément.)
    \item Combien de temps la lumière de l'éclair a-t-elle mis pour atteindre Léa ? Justifiez l'approximation faite en a).
    \item Le grondement du tonnerre a une fréquence de \qty{\v{f}}{Hz}. Déterminez sa longueur d'onde dans l'air.
\end{enumerate}
